\newpage
\vspace{-2.5cm}
\chapter*{\zihao{-3}\heiti{ABSTRACT}}

\vspace{-0.5cm}


\setlength{\baselineskip}{17pt}
With the continuous improvement of computer hardware performance and the enhancement of graphics processing capabilities, graphics simulation has been widely applied in many fields such as game development and film production. High-precision interactive simulation between fluids and elastoplastic solids has become a research hotspot. Although the Material Point Method (MPM) has achieved remarkable results in this scenario, there is still room for optimization from both the algorithmic level and the engineering perspective. For example, numerical dissipation occurs when simulating object collisions and fluid-solid coupling, blank gaps appear at the collision boundaries, and there is a lack of physical simulation software based on a node system in China.

Aiming at the above problems, this paper conducts in-depth research, and the main contributions are as follows: Firstly, in response to the limitation of numerical viscosity of the traditional MPM on the unified background grid, by modifying the traditional MPM process, a Decomposition-Compatible Affine Particle-In-Cell (DC-APIC) integrator is introduced. The free-slip boundary condition is enforced in the tangential direction, and the separation condition is achieved in the normal direction. Moreover, a segmentation method based on the phase-field gradient is designed. Secondly, to address the problem of white gaps generated when the traditional MPM deals with object collisions, an adaptive grid algorithm that incorporates the DC-APIC integrator is proposed. This algorithm can effectively retain collision details at different scales without incurring excessive additional costs. Thirdly, a physical simulation software based on the node system is designed and implemented. It implements the improved Material Point Method algorithm within the software framework. Through multiple challenging large-scale experimental scenarios, the robustness and practicality of the method in this paper in generating realistic visual effects are verified. This paper successfully provides a simulation method that can correctly simulate the proper separation of the interface in the two-way coupling between fluids and solids, as well as a set of node-based physical simulation software, which provides strong support for the development of related fields. 



%\hspace{-0.5cm}
{\sihao{\textbf{Keywords:}}} \textit{Computer Graphics; Physics-based Simulation; Material Point Method; Solid-fluid Coupling; Adaptive Methods}


%\clearpage
%\phantom{s}
%\clearpage
\addcontentsline{toc}{chapter}{ABSTRACT}